\documentclass[12pt,a4paper,french]{article}
\usepackage[utf8]{inputenc}
\usepackage{amsmath,amssymb,lmodern,babel,geometry}
\usepackage{pgf,tikz}
\usepackage{tkz-tab}
\usepackage{fancyhdr}
\usepackage{qrcode}
\usepackage{enumitem}
\usepackage[useregional]{datetime2}
\geometry{top=2.8cm, left=2cm, right=2cm, bottom=2cm}
\pagestyle{fancy}
\fancyhead[L]{Entraînement : Entraînement de calcul}
\fancyhead[R]{3e turquoise}
\fancyfoot[C]{\tiny{Date de création : \DTMnow}}
\renewcommand{\labelenumi}{\arabic{enumi})}
\setlist[enumerate]{topsep=0.2em,itemsep=0.2em,parsep=0pt,partopsep=0pt}
\setlength{\headheight}{15pt}
\begin{document}
\begin{center}
\textbf{Sujet numéro 1}
\\(à rendre avec la copie)
\end{center}


\noindent\textbf{Exercice 1} -- Fonctions de référence.
\begin{enumerate}
\item $f(x) = \dfrac{1}{x}. \text{ Comparer } f(-4) \text{ et } f(5)$
\item $f(x) = x^3. \text{ Comparer } f(3) \text{ et } f(1)$
\item $f(x) = \dfrac{1}{x}. \text{ Résoudre } f(x) = \dfrac{1}{2}$
\item $f(x) = \dfrac{1}{x}. \text{ Résoudre } f(x) = - \dfrac{3}{2}$
\item $f(x) = \sqrt{x}. \text{ Comparer } f(1) \text{ et } f(5)$
\end{enumerate}
\bigskip

\noindent\textbf{Exercice 2} -- Intervalles et valeur absolue.
\begin{enumerate}
\item $\text{Le nombre } -4 \text{ appartient-il à } \left[ -4 \,;\, 2 \right[ \text{ ?}$
\item $\text{Le nombre } 7 \text{ appartient-il à } \left] 6 \,;\, 12 \right] \text{ ?}$
\item $|x - (-5)| \leqslant 3$
\item $|x - (1)| < 3$
\item $|x - (4)| < 4$
\end{enumerate}
\bigskip

\vfill
\begin{flushright}
\qrset{height=1.5cm}
\qrset{level=M}
\qrcode{Ex1: < | > | x = 2 | x = - (2)/(3) | <  //  Ex2: Oui | Oui | [ -8 ; -2 ] | ] -2 ; 4 [ | ] 0 ; 8 [}
\end{flushright}
\newpage
\begin{center}
\textbf{Corrigé numéro 1}\\
\end{center}

\noindent\textbf{Exercice 1}\par
\smallskip
\begin{enumerate}
\item $f(-4) = - \dfrac{1}{4}$ et $f(5) = \dfrac{1}{5}$ donc $f(-4) < f(5)$
\item $f(3) = 27$ et $f(1) = 1$ donc $f(3) > f(1)$
\item $\dfrac{1}{x} = \dfrac{1}{2} \Longleftrightarrow x = 2$
\item $\dfrac{1}{x} = - \dfrac{3}{2} \Longleftrightarrow x = - \dfrac{2}{3}$
\item $\sqrt{x}$ est croissante sur $[0;+\infty[$. Comme $1 < 5$, $f(1) < f(5)$
\end{enumerate}
\medskip

\noindent\textbf{Exercice 2}\par
\smallskip
\begin{enumerate}
\item $-4$ $\in$ $\left[ -4 \,;\, 2 \right[$
\item $7$ $\in$ $\left] 6 \,;\, 12 \right]$
\item $|x - (-5)| \leqslant 3$

si et seulement si $-5 - 3 \leq x \leq -5 + 3$

si et seulement si $x \in \left[ -8 \,;\, -2 \right]$
\item $|x - (1)| < 3$

si et seulement si $1 - 3 < x < 1 + 3$

si et seulement si $x \in \left] -2 \,;\, 4 \right[$
\item $|x - (4)| < 4$

si et seulement si $4 - 4 < x < 4 + 4$

si et seulement si $x \in \left] 0 \,;\, 8 \right[$
\end{enumerate}
\medskip

\newpage
\begin{center}
\textbf{Sujet numéro 2}
\\(à rendre avec la copie)
\end{center}


\noindent\textbf{Exercice 1} -- Fonctions de référence.
\begin{enumerate}
\item $f(x) = \dfrac{1}{x}. \text{ Comparer } f(6) \text{ et } f(3)$
\item $f(x) = \sqrt{x}. \text{ Comparer } f(5) \text{ et } f(1)$
\item $f(x) = \dfrac{1}{x}. \text{ Comparer } f(-6) \text{ et } f(-2)$
\item $f(x) = x^3. \text{ Comparer } f(5) \text{ et } f(-3)$
\item $f(x) = \dfrac{1}{x}. \text{ Comparer } f(5) \text{ et } f(-1)$
\end{enumerate}
\bigskip

\noindent\textbf{Exercice 2} -- Intervalles et valeur absolue.
\begin{enumerate}
\item $\text{Le nombre } 15 \text{ appartient-il à } \left] 8 \,;\, 14 \right[ \text{ ?}$
\item $\text{Le nombre } 3 \text{ appartient-il à } \left[ 4 \,;\, 5 \right[ \text{ ?}$
\item $|x - (-4)| \leqslant 4$
\item $|x - (-7)| \leqslant 3$
\item $|x - (-8)| < 4$
\end{enumerate}
\bigskip

\vfill
\begin{flushright}
\qrset{height=1.5cm}
\qrset{level=M}
\qrcode{Ex1: < | > | > | > | >  //  Ex2: Non | Non | [ -8 ; 0 ] | [ -10 ; -4 ] | ] -12 ; -4 [}
\end{flushright}
\newpage
\begin{center}
\textbf{Corrigé numéro 2}\\
\end{center}

\noindent\textbf{Exercice 1}\par
\smallskip
\begin{enumerate}
\item $f(6) = \dfrac{1}{6}$ et $f(3) = \dfrac{1}{3}$ donc $f(6) < f(3)$
\item $\sqrt{x}$ est croissante sur $[0;+\infty[$. Comme $5 > 1$, $f(5) > f(1)$
\item $f(-6) = - \dfrac{1}{6}$ et $f(-2) = - \dfrac{1}{2}$ donc $f(-6) > f(-2)$
\item $f(5) = 125$ et $f(-3) = -27$ donc $f(5) > f(-3)$
\item $f(5) = \dfrac{1}{5}$ et $f(-1) = -1$ donc $f(5) > f(-1)$
\end{enumerate}
\medskip

\noindent\textbf{Exercice 2}\par
\smallskip
\begin{enumerate}
\item $15$ $\notin$ $\left] 8 \,;\, 14 \right[$
\item $3$ $\notin$ $\left[ 4 \,;\, 5 \right[$
\item $|x - (-4)| \leqslant 4$

si et seulement si $-4 - 4 \leq x \leq -4 + 4$

si et seulement si $x \in \left[ -8 \,;\, 0 \right]$
\item $|x - (-7)| \leqslant 3$

si et seulement si $-7 - 3 \leq x \leq -7 + 3$

si et seulement si $x \in \left[ -10 \,;\, -4 \right]$
\item $|x - (-8)| < 4$

si et seulement si $-8 - 4 < x < -8 + 4$

si et seulement si $x \in \left] -12 \,;\, -4 \right[$
\end{enumerate}
\medskip

\newpage
\end{document}
